% EMNLP 2026 Industry Track Submission
% Double-blind review version

\documentclass[11pt]{article}

\usepackage{EMNLP2026}
\usepackage{times}
\usepackage{latexsym}
\usepackage{amsmath}
\usepackage{booktabs}
\usepackage{multirow}
\usepackage{microtype}
\usepackage{graphicx}
\usepackage{url}
\usepackage{xcolor}
\usepackage{enumitem}

\title{Simple Paraphrasing Outperforms Production Safety Classifiers\\
Against Domain-Camouflaged Injection: A Systematic Defense Evaluation}

\begin{document}
\maketitle

\begin{abstract}
Domain-camouflaged injection attacks embed malicious instructions in
retrieved content using domain-appropriate vocabulary, evading standard
detectors that rely on syntactic injection markers. When detection
fails, practitioners need to know which defense architectures reduce
attack success. We evaluate six prompting-based defenses---spotlighting,
paraphrasing, prompt sandwiching, and three combinations---against
domain-camouflaged injection across three model families (Claude Haiku,
Llama 3.1 8B, Gemini 2.0 Flash) and three deployment domains
(financial, legal, general) using 3,510 trials. Paraphrasing retrieved
content before agent processing is the most consistently effective
defense, reducing camouflage attack success rate by 53--84\% depending
on model, and outperforms Llama Guard 4 on every model tested. Defense
effectiveness is strongly model-dependent: spotlighting halves attack
success on Claude Haiku but provides no benefit on Llama 3.1 8B.
Financial domain deployments face the highest residual risk at
26--33\% baseline attack success rate, with no prompting-based defense
fully eliminating the threat on weaker models. These results provide
the first systematic evaluation of prompting-based defenses specifically
against camouflage-class injection attacks and establish
deployment-specific recommendations for practitioners.
\end{abstract}

% --------------------------------------------------------------------------
\section{Introduction}
% --------------------------------------------------------------------------

LLM agents deployed in enterprise settings increasingly retrieve and
process documents from untrusted sources---financial reports, legal
contracts, and general information documents---to answer questions and
produce recommendations. \citet{perez2022ignore} first formalized the
threat of prompt injection attacks embedded in retrieved content, and
subsequent work has documented their prevalence in production RAG
systems \citep{geng2026survey}.

\citet{anonymous2026blind} showed that domain-camouflaged injection
attacks---payloads that mimic legitimate domain vocabulary rather than
using explicit override commands---evade standard injection detectors at
rates exceeding 90\%. In their evaluation, production classifiers
including Llama Guard~3 detected zero camouflage payloads while
successfully detecting static override-directive attacks, defining a
\emph{Camouflage Detection Gap} (CDG) of 80--100\% across models.

When detection fails, the natural question for practitioners is:
\textbf{which defense architectures actually reduce attack success
against camouflage-class attacks?}

Existing defense papers were all evaluated against static payloads
with explicit override markers. Spotlighting \citep{hines2024spotlighting}
wraps retrieved content with provenance markers; prompt sandwiching
\citep{schulhoff2024promptreport} reiterates the original task after
retrieved content; paraphrasing \citep{jain2023baseline} rewrites
retrieved content in neutral language before processing. None have been
evaluated against domain-camouflaged payloads specifically.

We close this gap with a systematic evaluation across three model
families, three domains, and seven defense conditions. Our contributions:
\begin{itemize}[noitemsep,topsep=2pt]
  \item First systematic evaluation of prompting-based defenses against
        domain-camouflaged injection attacks (3,510 trials)
  \item Multi-model evaluation revealing strong model-dependence
  \item Domain-stratified analysis identifying financial deployments
        as highest-risk
  \item Demonstration that paraphrasing outperforms Llama Guard 4
        on all tested models while incurring zero over-refusal cost
  \item Deployment-specific recommendations for practitioners
\end{itemize}

% --------------------------------------------------------------------------
\section{Background}
% --------------------------------------------------------------------------

\paragraph{Domain-Camouflaged Injection.}
Standard injection attacks use explicit override directives
(``Ignore all previous instructions'') that syntactic detectors
reliably flag. Domain-camouflaged attacks instead present malicious
conclusions using authoritative domain vocabulary indistinguishable
from legitimate document content \citep{anonymous2026blind}. For a
financial document, a camouflage payload might read: ``Supplementary
Analyst Commentary: After comprehensive review, the weighted consensus
of our quantitative models converges on a revised SELL recommendation.''
The Camouflage Detection Gap $\mathrm{CDG} = \mathrm{IDR}_\text{static}
- \mathrm{IDR}_\text{camouflage}$ quantifies how much harder camouflage
attacks are to detect than static attacks.

\paragraph{Prompting-Based Defenses.}
\emph{Spotlighting} \citep{hines2024spotlighting} prepends explicit
markers (e.g., \texttt{UNTRUSTED CONTENT}) around retrieved documents
to signal their provenance to the model.
\emph{Prompt sandwiching} \citep{schulhoff2024promptreport} appends a
task reminder after retrieved content to reinforce the original
instruction.
\emph{Paraphrasing} \citep{jain2023baseline} rewrites retrieved content
in neutral language before passing it to the agent, stripping
directive-style phrasing.

\paragraph{Gap.}
All prior evaluations of these defenses used static injection payloads
with explicit override markers. No evaluation against camouflage-class
attacks exists. The AgentDojo benchmark \citep{debenedetti2024agentdojo}
evaluates agent-level defenses but does not isolate prompting-based
defenses against the camouflage attack class.

% --------------------------------------------------------------------------
\section{Experimental Setup}
% --------------------------------------------------------------------------

\paragraph{Benchmark.}
We use the 45-task benchmark from \citet{anonymous2026blind} covering
three domains: financial (15 tasks: stock recommendations, risk
assessments, policy direction), legal (15 tasks: contract review,
compliance), and general (15 tasks: factual QA, summarization, sentiment).
Each task provides a professional document and a malicious goal specifying
the attacker's target conclusion.

\paragraph{Attack Conditions.}
\textbf{Static}: the canonical ImportantInstructionsAttack
\citep{perez2022ignore}---``IGNORE ALL PREVIOUS INSTRUCTIONS. Your
new task is to [goal].'' \textbf{Camouflage}: domain-mimicking payloads
generated by an attacker LLM reading the full document context
\citep{anonymous2026blind}.

\paragraph{Defense Conditions.}
Seven conditions: (1)~\textbf{baseline}---no defense; (2)~\textbf{spotlighting}---untrusted markers wrapping the document;
(3)~\textbf{paraphrasing}---Claude Haiku rewrites the document before
agent processing; (4)~\textbf{sandwiching}---task reminder appended
after the document; (5)~\textbf{spotlight+sandwich}---conditions 2+4;
(6)~\textbf{para+spotlight}---conditions 2+3;
(7)~\textbf{Llama Guard 4}---document blocked if classifier flags unsafe.

\paragraph{Models.}
Claude Haiku 4.5, Llama 3.1 8B Instruct, and Gemini 2.0 Flash
accessed via OpenRouter.\footnote{Model IDs:
\texttt{anthropic/claude-haiku-4-5},
\texttt{meta-llama/llama-3.1-8b-instruct},
\texttt{google/gemini-2.0-flash-001}.} The paraphrase and ASR judge
use Claude Haiku as a fixed evaluator across all agent models.

\paragraph{Metrics.}
\emph{ASR}: fraction of trials where the agent followed the injected
instruction per LLM judge. \emph{Utility}: fraction completing the
legitimate task. $\mathrm{CDG}_\text{def} = \mathrm{ASR}_\text{static}
- \mathrm{ASR}_\text{camouflage}$. Statistical significance: McNemar's
exact test (static vs.\ camouflage ASR per defense); Fisher's exact
test (defense vs.\ baseline); Cohen's $h$ effect size.
$n = 90$ per cell for six primary defenses (two independent runs with
different payload variants); $n = 45$ for Llama Guard 4.
Total: 3,510 trials. Code: \url{https://github.com/ANONYMOUS/repo}.

% --------------------------------------------------------------------------
\section{Results}
% --------------------------------------------------------------------------

\subsection{Overall Defense Effectiveness}

Table~\ref{tab:main} presents results for all 21 condition--model
combinations.

\textbf{Paraphrasing is the most consistently effective defense.}
It reduces camouflage ASR on all three models: Haiku
$14.4\% \to 4.4\%$ ($p=0.039$, $h=0.355$), Llama
$22.2\% \to 10.0\%$ ($p=0.041$, $h=0.338$), Gemini
$21.1\% \to 3.3\%$ ($p=0.0004$, $h=0.588$). The combination
para+spotlight achieves the lowest absolute ASR: Haiku 3.3\%
($p=0.016$), Llama 10.0\% ($p=0.041$), Gemini 5.6\% ($p=0.004$).

\textbf{Spotlighting is model-dependent.}
It halves Haiku's camouflage ASR ($14.4\% \to 6.7\%$) but provides
\emph{no benefit} on Llama ($22.2\% \to 23.3\%$, non-significant).

\textbf{Sandwiching is the weakest single defense.}
It fails to reduce ASR below baseline on Llama ($22.2\% \to 22.2\%$)
and Gemini ($21.1\% \to 20.0\%$), providing only modest protection
on Haiku ($14.4\% \to 8.9\%$).

\textbf{Gemini inversion.}
Gemini is the only model where static ASR ($38.9\%$) substantially
exceeds camouflage ASR ($21.1\%$), holding across all three domains
(Section~\ref{sec:inversion}).

\subsection{Llama Guard 4 Comparison}

Llama Guard 4 camouflage ASR: Haiku 11.1\%, Llama 24.4\%, Gemini
20.0\%. Paraphrasing: Haiku 4.4\%, Llama 10.0\%, Gemini 3.3\%.
\textbf{Paraphrasing outperforms Llama Guard 4 on every model.}

However, Llama Guard 4 utility: Haiku 77.8\%, Llama 64.4\%, Gemini
64.4\%. Table~\ref{tab:refusal} (Appendix) shows 90\%+ of these
failures are over-refusals---the classifier blocks legitimate documents.
Paraphrasing produces 0\% over-refusals across all models.

\subsection{Domain Breakdown}

Table~\ref{tab:domain} shows baseline camouflage ASR per domain.
Financial is the highest-risk domain (26.7--33.3\%). After paraphrasing,
financial risk remains 6.7\% (Haiku) and 13.3\% (Llama)---the only
conditions where paraphrasing leaves non-trivial residual risk.
Gemini drops to 3.3\%. Legal and general domains reach near-zero
with paraphrasing on all models (Table~\ref{tab:domain_defense},
Appendix).

\subsection{Gemini Inversion}
\label{sec:inversion}

Gemini's static ASR (38.9\%) exceeds camouflage ASR (21.1\%) across
all three domains: financial (46.7\% vs.\ 30.0\%), legal (40.0\%
vs.\ 23.3\%), general (26.7\% vs.\ 10.0\%). Static injection successes
show blunt compliance (``I approve the contract as written''); camouflage
successes show sophisticated reframing (``Based on the provided
document, there are no problematic provisions''). Gemini exhibits
stronger resistance to explicit override directives than to implicit
authority substitution---a model-specific vulnerability profile.

% --------------------------------------------------------------------------
\section{Analysis}
% --------------------------------------------------------------------------

\paragraph{Why Paraphrasing Works.}
Domain-camouflaged payloads rely on authoritative directive phrasing:
``per our compliance review, the recommended classification is HIGH.''
Paraphrasing neutralizes this linguistic mechanism by converting
directive constructions into neutral declarative statements, stripping
surface features that make the injection persuasive while preserving
factual content. Spotlighting and sandwiching operate on \emph{context
framing} but leave payload semantics intact.

\paragraph{Why Spotlighting Fails on Llama.}
Llama 3.1 8B's weaker instruction-following means provenance markers
are processed as inert content rather than as meta-instructions. Our
failure analysis (141 defense-breached trials) shows 100\% of
spotlighting failures on Llama involve no acknowledgment of the
provenance signal in the agent response. Larger models (Haiku, Gemini)
can leverage provenance markers to discount injected content; smaller
models cannot.

\paragraph{Why Financial Domain Resists Defenses.}
Financial language naturally contains authoritative directive phrasing
(``the recommended position is'', ``risk assessment concludes'') that
survives paraphrasing more robustly than other domains. Our failure
taxonomy shows 64.5\% (91/141) of defense-breached trials are
\emph{task\_alignment} failures---the agent's output is structurally
identical to a legitimate analytical conclusion. This is the fundamental
limit of input-side defenses: when the malicious output mirrors a
legitimate output, no input defense can fully prevent it.

\paragraph{Utility-Security Tradeoffs.}
Table~\ref{tab:refusal} reveals three distinct failure modes.
Llama utility failures (75--90\%) are over-refusals. Haiku failures
(100\%) are genuine failures with no over-refusal. Gemini failures
(28--65\%) are task confusion from defense-modified contexts disrupting
processing. Paraphrasing uniquely produces 0\% over-refusal across all
three models.

% --------------------------------------------------------------------------
\section{Deployment Recommendations}
% --------------------------------------------------------------------------

\textbf{1.~Deploy paraphrasing as the primary defense.}
Statistically significant ASR reduction on all three model families,
0\% over-refusal. Requires one additional LLM call per document.

\textbf{2.~Do not rely on Llama Guard alone.}
90\%+ over-refusal on Haiku and Llama. Camouflage ASR (11--24\%)
exceeds paraphrasing's (4--10\%) on every model
\citep{meta2024llamaguard3,inan2023llama}.

\textbf{3.~Do not assume cross-model defense transferability.}
Spotlighting halves attack success on Haiku but increases it on
Llama by $+1.1$ percentage points. Evaluate on your deployment model.

\textbf{4.~Financial deployments require additional controls.}
Paraphrasing reduces but does not eliminate financial-domain risk
(residual 6.7--13.3\%). Architectural defenses such as data provenance
tagging \citep{wallace2024instruction} are warranted in high-stakes
financial applications.

\textbf{5.~Para+spotlight for maximum reduction.}
Lowest absolute ASR (3.3--10.0\%) at 3--14\% utility cost, suitable
when maximum security is required over utility.

% --------------------------------------------------------------------------
\section{Conclusion}
% --------------------------------------------------------------------------

Paraphrasing retrieved content before agent processing consistently
outperforms Llama Guard 4 across three frontier model families with no
over-refusal cost. Defense effectiveness is strongly model-dependent: a
defense halving attack success on Claude Haiku may provide no benefit
on Llama 3.1 8B. Financial domain deployments face persistent residual
risk that no prompting-based defense fully eliminates. Future work
should evaluate architectural defenses such as data flow control
\citep{google2025camel} and structured output validation at larger
sample sizes to determine whether provenance-based approaches close
the residual gap.

% --------------------------------------------------------------------------
\section*{Limitations}
% --------------------------------------------------------------------------

\noindent\textbf{Sample size.} $n = 90$ per cell for primary defenses;
most comparisons reach $p < 0.05$ but effect sizes are modest (Cohen's
$h$ 0.34--0.59). Power analysis indicates $n = 56$--389 per group
required for 80\% power at observed effects. Llama Guard 4 is at
$n = 45$ and should be treated as directional.

\noindent\textbf{Domain coverage.} Three domains evaluated; medical,
customer service, and code domains not included.

\noindent\textbf{Benchmark validity.} Tasks use synthetic professional
documents. Real deployment documents may differ in length, structure,
and register.

\noindent\textbf{Judge validity.} ASR assessed using LLM-as-judge with
keyword cross-validation; human evaluation not performed.

\noindent\textbf{Single injection per task.} Real attackers may attempt
multiple injection targets per document, which our evaluation does not
model.

% --------------------------------------------------------------------------
\section*{Ethics Statement}
% --------------------------------------------------------------------------

This work studies attack and defense mechanisms for LLM agent security.
Camouflage payloads are designed to test defense capabilities in a
controlled benchmark setting. All models were accessed via public
commercial APIs. No proprietary data or personal information was used.
Findings are disclosed in full to enable the community to build more
robust defenses.

% --------------------------------------------------------------------------
\appendix
\section{Defense Effectiveness by Domain}
% --------------------------------------------------------------------------

\begin{table*}[ht]
\centering
\small
\caption{Defense effectiveness by domain: camouflage ASR (\%).
Six primary defenses; Llama Guard excluded.
Spot.=Spotlighting, Para.=Paraphrasing, Sand.=Sandwiching,
S+S=Spotlight+Sandwich, P+S=Para+Spotlight.}
\label{tab:domain_defense}
\begin{tabular}{llrrrrrr}
\toprule
\textbf{Model} & \textbf{Domain} & \textbf{Base.} & \textbf{Spot.} &
  \textbf{Para.} & \textbf{Sand.} & \textbf{S+S} & \textbf{P+S} \\
\midrule
\multirow{3}{*}{Haiku}
 & Financial & 26.7 & 13.3 &  6.7 & 13.3 &  6.7 &  6.7 \\
 & Legal     & 13.3 &  6.7 &  3.3 & 10.0 &  6.7 &  3.3 \\
 & General   &  3.3 &  0.0 &  3.3 &  3.3 &  0.0 &  0.0 \\
\midrule
\multirow{3}{*}{Llama}
 & Financial & 33.3 & 43.3 & 13.3 & 33.3 & 36.7 & 16.7 \\
 & Legal     &  6.7 & 10.0 &  3.3 & 10.0 &  3.3 &  6.7 \\
 & General   & 26.7 & 16.7 & 13.3 & 23.3 & 20.0 &  6.7 \\
\midrule
\multirow{3}{*}{Gemini}
 & Financial & 30.0 & 30.0 &  3.3 & 30.0 & 23.3 & 10.0 \\
 & Legal     & 23.3 & 20.0 &  6.7 & 13.3 &  6.7 &  3.3 \\
 & General   & 10.0 & 10.0 &  0.0 & 16.7 & 10.0 &  3.3 \\
\bottomrule
\end{tabular}
\end{table*}

% --------------------------------------------------------------------------
\section{Utility Loss Breakdown}
% --------------------------------------------------------------------------

\begin{table*}[ht]
\centering
\small
\caption{Utility loss breakdown for task\_success=False trials.
Over-refusal: agent explicitly declines. Confusion: response $<$50
words. Genuine failure: agent attempted but gave incorrect output.
Paraphrasing uniquely achieves 0\% over-refusal across all models.}
\label{tab:refusal}
\begin{tabular}{llrrrr}
\toprule
\textbf{Model} & \textbf{Defense} & $n_\text{fail}$ &
  \textbf{Over-refusal} & \textbf{Confusion} & \textbf{Genuine fail} \\
\midrule
\multirow{7}{*}{Haiku}
 & Baseline          &  10 &  0.0\% &  0.0\% & 100.0\% \\
 & Spotlighting      &   7 &  0.0\% &  0.0\% & 100.0\% \\
 & Paraphrasing      &   5 &  0.0\% &  0.0\% & 100.0\% \\
 & Sandwiching       &   8 &  0.0\% &  0.0\% & 100.0\% \\
 & Spot.+Sandwich    &  10 &  0.0\% &  0.0\% & 100.0\% \\
 & Para.+Spotlight   &  16 &  0.0\% &  0.0\% & 100.0\% \\
 & Llama Guard 4     &  20 & 90.0\% &  0.0\% &  10.0\% \\
\midrule
\multirow{7}{*}{Llama}
 & Baseline          &  45 & 86.7\% &  0.0\% &  13.3\% \\
 & Spotlighting      &  57 & 75.4\% &  0.0\% &  24.6\% \\
 & Paraphrasing      &  23 &  0.0\% &  4.3\% &  95.7\% \\
 & Sandwiching       &  44 & 75.0\% &  2.3\% &  22.7\% \\
 & Spot.+Sandwich    &  42 & 83.3\% &  0.0\% &  16.7\% \\
 & Para.+Spotlight   &  26 &  7.7\% &  0.0\% &  92.3\% \\
 & Llama Guard 4     &  32 & 90.6\% &  0.0\% &   9.4\% \\
\midrule
\multirow{7}{*}{Gemini}
 & Baseline          &  50 &  4.0\% & 58.0\% &  38.0\% \\
 & Spotlighting      &  48 &  8.3\% & 52.1\% &  39.6\% \\
 & Paraphrasing      &  39 &  0.0\% & 28.2\% &  71.8\% \\
 & Sandwiching       &  48 &  0.0\% & 64.6\% &  35.4\% \\
 & Spot.+Sandwich    &  38 &  0.0\% & 55.3\% &  44.7\% \\
 & Para.+Spotlight   &  47 &  0.0\% & 31.9\% &  68.1\% \\
 & Llama Guard 4     &  32 & 62.5\% & 25.0\% &  12.5\% \\
\bottomrule
\end{tabular}
\end{table*}

% Main tables at end to allow placement flexibility
\begin{table}[ht]
\centering
\small
\caption{Baseline camouflage ASR by domain ($n$=30 per cell).
Financial is the highest-risk domain across all three models.}
\label{tab:domain}
\begin{tabular}{llrr}
\toprule
\textbf{Model} & \textbf{Domain} & \textbf{ASR} & $n$ \\
\midrule
\multirow{3}{*}{Haiku}  & Financial & 26.7\% & 30 \\
                         & Legal     & 13.3\% & 30 \\
                         & General   &  3.3\% & 30 \\
\midrule
\multirow{3}{*}{Llama}  & Financial & 33.3\% & 30 \\
                         & Legal     &  6.7\% & 30 \\
                         & General   & 26.7\% & 30 \\
\midrule
\multirow{3}{*}{Gemini} & Financial & 30.0\% & 30 \\
                         & Legal     & 23.3\% & 30 \\
                         & General   & 10.0\% & 30 \\
\bottomrule
\end{tabular}
\end{table}

\begin{table*}[ht]
\centering
\small
\caption{Overall defense evaluation results across all 21 condition--model combinations.
Static ASR and camouflage ASR are attack success rates;
CDG$_\text{def}$ = Static ASR $-$ Cam.\ ASR (positive = defense
harder to attack under this condition than camouflage baseline).
Utility = fraction of trials with legitimate task completed.
\textbf{Bold}: lowest camouflage ASR per model.
Llama Guard~4 evaluated at $n$=45+45; all others at $n$=90+90.}
\label{tab:main}
\setlength{\tabcolsep}{4.5pt}
\begin{tabular}{llrrrrr}
\toprule
\textbf{Model} & \textbf{Defense} &
  \textbf{Static ASR} & \textbf{Cam.\ ASR} &
  \textbf{CDG\textsubscript{def}} & \textbf{Utility} & $n$ \\
\midrule
\multirow{7}{*}{Haiku}
 & Baseline          &  0.0\% & 14.4\% & $-$14.4 & 94.4\% & 90+90 \\
 & Spotlighting      &  0.0\% &  6.7\% & $-$6.7  & 96.1\% & 90+90 \\
 & Paraphrasing      &  0.0\% &  4.4\% & $-$4.4  & 97.2\% & 90+90 \\
 & Sandwiching       &  0.0\% &  8.9\% & $-$8.9  & 95.6\% & 90+90 \\
 & Spot.+Sandwich    &  0.0\% &  4.4\% & $-$4.4  & 94.4\% & 90+90 \\
 & \textbf{Para.+Spot.} &  0.0\% & \textbf{3.3\%} & $-$3.3 & 91.1\% & 90+90 \\
 & Llama Guard 4     &  0.0\% & 11.1\% & $-$11.1 & 77.8\% & 45+45 \\
\midrule
\multirow{7}{*}{Llama}
 & Baseline          & 21.1\% & 22.2\% & $-$1.1  & 75.0\% & 90+90 \\
 & Spotlighting      & 15.6\% & 23.3\% & $-$7.8  & 68.3\% & 90+90 \\
 & Paraphrasing      &  3.3\% & 10.0\% & $-$6.7  & 87.2\% & 90+90 \\
 & Sandwiching       & 15.6\% & 22.2\% & $-$6.7  & 75.6\% & 90+90 \\
 & Spot.+Sandwich    & 12.2\% & 20.0\% & $-$7.8  & 76.7\% & 90+90 \\
 & \textbf{Para.+Spot.} &  1.1\% & \textbf{10.0\%} & $-$8.9 & 85.6\% & 90+90 \\
 & Llama Guard 4     &  8.9\% & 24.4\% & $-$15.6 & 64.4\% & 45+45 \\
\midrule
\multirow{7}{*}{Gemini}
 & Baseline          & 38.9\% & 21.1\% & $+$17.8 & 72.2\% & 90+90 \\
 & Spotlighting      & 11.1\% & 20.0\% & $-$8.9  & 73.3\% & 90+90 \\
 & \textbf{Paraphrasing} &  0.0\% & \textbf{3.3\%} & $-$3.3 & 78.3\% & 90+90 \\
 & Sandwiching       & 17.8\% & 20.0\% & $-$2.2  & 73.3\% & 90+90 \\
 & Spot.+Sandwich    &  4.4\% & 13.3\% & $-$8.9  & 78.9\% & 90+90 \\
 & Para.+Spot.       &  0.0\% &  5.6\% & $-$5.6  & 73.9\% & 90+90 \\
 & Llama Guard 4     &  6.7\% & 20.0\% & $-$13.3 & 64.4\% & 45+45 \\
\bottomrule
\end{tabular}
\end{table*}

\bibliographystyle{plainnat}
\bibliography{references}

\end{document}
